\documentclass[11pt,a4paper]{article}
\usepackage[utf8]{inputenc}
\usepackage[T1]{fontenc}
\usepackage{geometry}
\usepackage{amsmath}
\usepackage{booktabs}
\usepackage{enumitem}
\usepackage{hyperref}
\usepackage{parskip}

\geometry{margin=1in}
\title{\textbf{Metrik AI: Building Energy Consumption Prediction}\\[0.5em]
\large A Real-World Project for Advanced Python for Data Science}
\author{DS-GA 1019 \,|\, Spring 2026 \,|\, Project Proposal\\[0.5em]
\small Team: to be submitted by Feb 27}
\date{}

\begin{document}
\maketitle

\begin{abstract}
We propose a project on \textbf{building energy consumption prediction and operational support} using the ASHRAE Great Energy Predictor III dataset (53.6 million hourly meter readings). We deliver three components: (1)~scalable \textbf{forecasting} (next-hour consumption, baseline vs.\ LightGBM), (2)~residual-based \textbf{anomaly detection} at scale, and (3)~\textbf{decision support} (ranked audit list). The pipeline is out-of-core, leakage-safe, and demonstrates advanced Python: chunked I/O, vectorization, multiprocessing, Numba, profiling, and a structured codebase with CLI, config, and tests---addressing energy waste, cost, and carbon emissions.
\end{abstract}

\section{The Real-World Problem}

\subsection{Who Has the Problem?}
Building owners, facility managers, utilities, and sustainability teams.

\subsection{What Is the Problem?}
Commercial buildings waste enormous amounts of energy because heating, cooling, and lighting are often run on fixed schedules or reactive rules instead of actual demand. This leads to:
\begin{itemize}[noitemsep]
    \item \textbf{High cost} --- Energy is one of the largest operating expenses for buildings.
    \item \textbf{Carbon emissions} --- Buildings account for a large share of global energy use and GHG emissions.
    \item \textbf{Comfort vs.\ efficiency trade-offs} --- Without good predictions, operators either over-cool/over-heat (waste) or under-deliver (complaints).
\end{itemize}

\subsection{Why Does Prediction Help?}
Accurate prediction of \textbf{next-hour or next-day energy use} enables:
\begin{itemize}[noitemsep]
    \item Pre-cool/pre-heat at the right times (demand response).
    \item Flagging anomalies (e.g., meter malfunction or equipment left on).
    \item Benchmarking buildings and prioritizing retrofits.
    \item Supporting sustainability and net-zero targets with data.
\end{itemize}

\textbf{Problem statement:} Predict building energy consumption so that operators can optimize usage, cut cost, and reduce emissions.

\section{The Data}

\textbf{Dataset:} ASHRAE --- Great Energy Predictor III (Kaggle).\\
\url{https://www.kaggle.com/competitions/ashrae-energy-prediction}

\begin{itemize}[noitemsep]
    \item \textbf{Scale:} $\sim$53.6 million rows of hourly meter readings.
    \item \textbf{Buildings:} 1,636 non-residential buildings; 3,053 energy meters.
    \item \textbf{Geography:} 19 sites across North America and Europe.
    \item \textbf{Time span:} 2 full years (2016--2017).
    \item \textbf{Meter types:} Whole-building electricity, heating/cooling water, steam, solar, water/irrigation.
    \item \textbf{Auxiliary data:} Building metadata (square footage, year built, primary use) and weather (air temperature).
\end{itemize}

The dataset is publicly available, requires no scraping or external APIs, and is large enough to justify scalable processing (chunked reads, Dask, vectorized feature engineering, optional Numba for heavy computations).

\section{The Solution}

\subsection{Framing: Forecasting + Anomaly Detection + Decision Support}
We deliver three components to maximize operational impact and demonstrate advanced Python at scale: (1)~\textbf{Forecasting} --- next-hour meter-level energy consumption using out-of-core, leakage-safe features and a baseline vs.\ LightGBM comparison. (2)~\textbf{Anomaly detection} --- residual-based scoring (actual minus predicted) to flag malfunctioning meters or abnormal consumption, implemented efficiently over large arrays. (3)~\textbf{Decision support} --- a ranked output (e.g., top meters or buildings to audit) so operators know where to act first.

\subsection{Data Processing Strategy (50M+ Rows)}
\textbf{Out-of-core by default.} The main train table is never fully loaded into memory. We use chunked reads (e.g., 1--2M rows per chunk) with pandas or Dask; merge each chunk with in-memory building metadata and weather; then either write features to disk (partitioned by site or time) or feed a downstream sampler. Peak memory target: under 4\,GB per process. In-memory: building metadata, weather (per site), and any per-site or per-meter subset used for a single model.

\subsection{Feature Engineering (Time-Aware, Leakage-Safe)}
Features use only past and current information. \textit{Time:} hour, day of week, month, is\_weekend, holiday (calendar). \textit{Lags:} same hour yesterday, same day last week. \textit{Rolling:} 24h and 168h rolling mean with an expanding or correctly aligned window so no future data leaks. \textit{Building:} square footage, primary use, year built. \textit{Weather:} air temperature (and optionally lagged). Implementation: vectorized NumPy/pandas per chunk; no Python loops over rows. Unit tests assert that no feature depends on future timestamps.

\subsection{Modeling Strategy}
\textbf{Baseline:} Per-meter mean or ``same hour yesterday'' predictor; time-based train/validation split (e.g., last 3--6 months held out). \textbf{Stronger model:} LightGBM (handles large data, categoricals, fast training). One global model with site/meter identifiers as categoricals, or per-site models trained in parallel. Target: continuous meter consumption; metric: RMSE or CV-RMSE. Success target: LightGBM achieves at least 20--40\% RMSE improvement over baseline on the same split.

\subsection{Scalability Strategy}
\textbf{Partitioning:} by site (19 sites) or by time chunks. \textbf{Parallelization:} multiprocessing or concurrent.futures --- e.g., one worker per site for feature build or training, or parallel chunks for feature construction. No shared mutable state; each worker reads its partition and writes outputs. Benchmark: report runtime and speedup (e.g., 1 vs.\ 4 workers) on a fixed subset (e.g., 5M rows).

\subsection{Acceleration Strategy}
\textbf{Numba JIT} on the hottest path: e.g., anomaly scoring over millions of residuals (score per meter or per timestep) or a custom rolling aggregation not easily expressed in pandas. Use \texttt{@jit(nopython=True)}; compare runtime vs.\ pure Python loop. If a bottleneck is in pandas/NumPy, vectorization is preferred before JIT. No Cython unless profiling shows a clear need.

\subsection{Profiling and Benchmarking}
\textbf{Tools:} \texttt{cProfile} for CPU hotspots, \texttt{memory\_profiler} for peak memory. \textbf{Metrics:} runtime (s), peak memory (GB), speedup (optimized/baseline). \textbf{Comparisons:} (1)~data load: naive full read (or OOM) vs.\ chunked; (2)~feature build: single-thread vs.\ vectorized and/or parallel; (3)~anomaly: Python loop vs.\ Numba or vectorized. One benchmark table in the report with realistic numbers; at least one bottleneck identified and improved.

\subsection{Engineering Strategy}
\textbf{Structure:} \texttt{src/} (load, features, model, anomaly, decision\_support), \texttt{scripts/}, \texttt{config/}, \texttt{tests/}. \textbf{CLI:} one entry point, e.g.\ \texttt{python -m src.cli run --config config.yaml}. \textbf{Config:} YAML/TOML for paths, \texttt{chunk\_size}, \texttt{seed}, model params. \textbf{Logging:} standard library or structlog; INFO level. \textbf{Typing:} type hints on public functions. \textbf{Tests:} pytest; at least feature-leakage check and one small end-to-end run. \textbf{Reproducibility:} \texttt{requirements.txt}, fixed seed in config, README with run instructions.

\subsection{Benchmark Table (Targets)}
\begin{table}[h]
\centering
\small
\begin{tabular}{@{}lccccc@{}}
\toprule
\textbf{Component} & \textbf{Baseline} & \textbf{Optimized} & \textbf{Runtime} & \textbf{Peak mem.} & \textbf{Success target} \\
\midrule
Data load & Full read / OOM & Chunked & --- & $<$4\,GB & No OOM \\
Feature build (5M rows) & Single-thread & Vectorized $\pm$ parallel & --- & --- & $\geq 2\times$ speedup \\
Forecasting & Mean predictor & LightGBM & --- & --- & 20--40\% RMSE gain \\
Anomaly (1M rows) & Python loop & Numba / vectorized & --- & --- & $\geq 5\times$ speedup \\
\bottomrule
\end{tabular}
\caption{Benchmark template: fill with measured runtimes and memory.}
\end{table}

\subsection{Staged Development (Feasibility)}
Phase 1: Single-site prototype (chunked I/O, features, baselines + LightGBM). Phase 2: Profiling and optimization targets. Phase 3: Anomaly scoring and audit ranking; benchmark table populated. Phase 4: Parallel scaling; optional full 53M run. Phase 5: Final report (4 pages), code, presentation. Subset by site or rows if needed to meet timeline.

\subsection{Deliverables}
\begin{itemize}[noitemsep]
    \item Reproducible pipeline (scripts + config) producing forecasts, anomaly scores, and a decision-support list.
    \item Benchmark table with baseline vs.\ optimized runtimes and memory; at least one documented speedup.
    \item Final report ($\leq$4 pages): problem, data, solution (three components), benchmarks, results, limitations.
    \item Code: typed, tested (leakage check + one integration test), CLI, README.
\end{itemize}

\section{Why This Is a Strong Choice}

\begin{table}[h]
\centering
\begin{tabular}{@{}ll@{}}
\toprule
\textbf{Criterion} & \textbf{How it fits} \\
\midrule
Real-world & Energy waste and carbon from buildings are a major, well-documented issue. \\
Stakeholders & Building operators, utilities, ESG/sustainability teams. \\
Data & One public Kaggle dataset, 53M rows; no scraping or APIs required. \\
Feasibility & Subset by site or meter count and still show a complete solution; full scale optional. \\
Advanced Python & Chunked I/O, Dask, vectorization, multiprocessing, Numba, profiling. \\
Course alignment & Data science + optimization at scale; clear story for proposal and final presentation. \\
\bottomrule
\end{tabular}
\caption{Project fit to course and real-world criteria.}
\end{table}

\section{Alternative Project Ideas}

\subsection{Alternative 1: Store Sales / Demand Forecasting (Retail)}
\textbf{Problem:} A grocery retailer (e.g., Corporación Favorita) needs to predict unit sales per product per store to reduce stockouts and food waste. \\
\textbf{Data:} Kaggle ``Store Sales - Time Series Forecasting'' (Favorita). \\
\textbf{Solution:} Time series (or hierarchical) forecasting with date, store, product, promotions, holidays; vectorized features and optional parallel runs per store or product family. \\
\textbf{Impact:} Less waste, better product availability.

\subsection{Alternative 2: Air Quality (PM2.5) Prediction}
\textbf{Problem:} Cities and health agencies need daily or sub-daily PM2.5 estimates to issue advisories and protect vulnerable populations. \\
\textbf{Data:} NASA Airathon (Taipei, Delhi, LA) or Kaggle PM2.5 datasets. \\
\textbf{Solution:} Predict PM2.5 from satellite-derived AOD, weather, and optionally topography; spatial--temporal modeling; optional Numba for custom spatial aggregates. \\
\textbf{Impact:} Public health, outdoor activity guidance, policy.

\section{Summary}

\textbf{Recommended project:} Building energy consumption prediction using the ASHRAE dataset---a clear real-world problem (cost and carbon), a 53M-row dataset, and a concrete solution (forecasting pipeline with scalable, advanced Python). \\
\textbf{Alternatives:} Retail demand forecasting (Favorita) or air quality (PM2.5) if a different domain is preferred; both are real-world, have public data, and admit clear solutions.

\end{document}
